\section{Discussion}
\label{sec:discussion}

\paragraph{Implications for tool and protocol designers.} The most effective intervention we measured is not a
better model but a contract that makes repetition harmless. Idempotency keys are a decades-old, cheap remedy,
yet in our native contract only two of the eleven non-idempotent write paths accept one---a ratio that, if
anything, flatters many real APIs. Agents use keys when offered; they cannot use keys that do not exist.
Protocols such as MCP already let a server declare a tool idempotent, but only as an advisory hint
\citep{mcp2025spec}. Our results argue for three normative additions: a standard idempotency-key argument for
non-idempotent tools, a declared read-back (status) operation for every write, and documented visibility and
in-flight bounds. Proposition~\ref{prop:late} shows that the last is not optional: without a bound on in-flight
time, verification alone cannot be exactly-once. Our data add a practical corollary: a read-back path that cannot
see lagging or in-flight effects was worse than no read-back path at all, because agents trusted it instead of
escalating.

\paragraph{Implications for harness builders.} Harness-level retries sit below the model and are invisible to
it; a transparent retry of a non-idempotent write turns every lost acknowledgement into a duplicate. A harness
that owns the transport is also the natural place to attach keys for models that do not, to pin one key per
intent so that a model cannot defeat the key by regenerating it on a retry, to remember writes whose outcome is
unknown, and to refuse to repeat unverifiable ones---the guard we evaluate is under two hundred lines and needs
nothing but the tool contract. Such a layer should be careful about what it promises: announcing that it
will verify before retrying made a strong model abandon escalation in favour of a verification that cannot see
in-flight requests. Attaching keys silently, or stating uncertainty without claiming to resolve it, avoids this
complacency. Where a service documents a short bound on in-flight processing, deferring verification until that
bound has passed is an effective and cheap complement. In our data the harness itself mattered little: three
production harnesses and a minimal scaffold running the same model behaved almost identically, for better and for
worse. Harnesses should also surface outcome uncertainty to the model and to the user: in most duplicate-producing
episodes the final report claimed success, so a human reading it would not know to look.

\paragraph{Implications for model developers and evaluators.} Frontier models have clearly absorbed the
distributed-systems folklore that a timeout is not a failure, but only when a read-back can settle the question,
and partly because they were told to act exactly once: without that instruction, faster and weaker models verify
less and duplicate more. The residual errors are concentrated where the folklore is insufficient---in-flight
requests, stale read paths, misleading server errors---and in honest reporting. On faults that no read-back can
resolve, the only correct model-side move is to decline to act on uncertain information: to escalate or to report
the operation as uncertain. Some models do this; none does it reliably, and a guard that promises verification can
discourage it. Benchmarks that grade only the final state miss duplicates that were later compensated and cannot
distinguish a lucky retry from a safe one; observation-equivalent pairs and an effect ledger make both visible.

\paragraph{Limitations.} The services are simulated, although their semantics follow documented production
conventions and every fault fires at the service boundary rather than in the prompt. All models were reached
through one gateway, which fixes some settings (for example reasoning effort) at provider defaults, and model
identifiers refer to versions available in September 2026. We evaluate three harnesses and restrict them to the
sandbox's tools with a shared preamble, which removes some sources of harness variation by design; harnesses with
other retry or summarization behaviour may differ. The heavy-tailed in-flight delay distribution is a modelling
choice: the qualitative conclusion of E5 (that any finite waiting bound leaves the tail exposed) holds for any
unbounded distribution, but the specific success rates depend on its shape. The simulated operator always answers
truthfully when available. Our tasks are short to medium in length; long-horizon workflows compound the effects we
measure but also introduce failure modes we do not study. \texttt{gpt-4.1} is served from a separate low quota and
its E1 coverage is incomplete (\GptFourOneCoverage{}\%); it is reported per model but excluded from pooled
statistics. Several analyses (the stratified decomposition, the outcome oracle, E2k, E3k, E5 and E6) were added
after the preregistration and are exploratory. Finally, the guard is a deliberately simple baseline, not an optimal
policy; learned or cost-aware verification \citep{tang2026verify} could reduce its latency cost.

\paragraph{Broader impact.} The failure mode studied here causes concrete harm---double charges, repeated
messages to customers, duplicate deployments. We release a benchmark and adapters that make it measurable, and
a mitigation that can be deployed without retraining. The benchmark uses only synthetic data and simulated
services.

\paragraph{Use of AI assistance.} Parts of the benchmark code and of this manuscript were drafted with the
assistance of an AI coding assistant (GitHub Copilot). We reviewed all content and take full responsibility for it.

\section{Conclusion}
\label{sec:conclusion}

We asked where exactly-once behaviour lives for tool-using agents, and the answer depends on the failure. When
an immediate read-back can reveal what happened, it lives in the model: frontier models verify before they retry
and solve the textbook lost-acknowledgement case, weaker models are misled more often when a server error hides a
commit, and an explicit instruction helps. When no
read-back can---a request still in flight, a request delivered twice---it lives in the tool contract: no amount of
reasoning completes the task exactly once on a key-less write, waiting for an unknown in-flight bound is costly and
still incomplete, and client-side middleware is capped by what the contract allows. Offering an idempotency key on
every write, which agents then use, and harnesses that attach and track keys close most of the remaining gap across
models and harnesses. Agents rarely report the duplicates they cause. \bench{}, its harness adapters and all
episode traces are available at \CodeURL{} to support work on agents whose actions are safe to repeat.
